\chapter{关键功能实现}
\section{预约创建}
\subsection{输入校验}
创建接口首先解析日期，并检查起止分钟数是否为整数、是否位于开放时段内，以及结束时刻是否大于开始时刻。无效输入直接返回结果码，不启动写事务。进入事务后，再核对用户和座位是否存在，从而区分输入错误与未知资源。

\subsection{冲突检查与记录写入}
预约服务使用参数化查询传递用户编号、日期和时间值。座位冲突查询与用户冲突查询共用式\eqref{eq:overlap}的判定条件，分别按座位编号和用户编号限制对象。检测到冲突时返回相应结果码并回滚；没有冲突时生成预约记录，提交事务后返回记录编号。附录 B 给出了该接口的核心代码。

正常预约的处理顺序为输入检查、开启事务、对象检查、座位检查、用户检查、插入和提交。业务错误与数据库异常分别处理：业务错误转换为明确结果码，数据库异常在回滚后继续抛出，方便调用方定位运行故障。

\section{预约取消}
取消功能执行一条带条件的更新语句。更新条件包含预约编号、当前用户编号及有效状态，成功修改一条记录时返回取消成功。条件不满足时，受影响的记录数为零。后续冲突检查只读取有效预约，因此成功取消后，原时段可以再次被其他用户预约。

\section{查询与运行结果}
统计接口按日期读取有效预约，分别计算记录数量和预约分钟数之和。原型同时保留取消记录，避免将状态变化等同于记录删除。调用方根据结果码展示预约成功、座位冲突、用户冲突、输入无效或取消失败等结果。

用户编号由测试程序直接传入，代表已确定的调用身份。实际应用接入校园账户后，应由认证会话提供这一编号，并在服务端执行同样的归属校验。预约服务因此可以在保持业务规则一致的前提下接入网页或移动端入口。
